% The layout of the LaTeX build, which `transpile.py` puts before the body
% it writes from `main.typ`. The journal's class replaces `article` at
% acceptance.
\documentclass[a4paper,10pt]{article}
\usepackage[T1]{fontenc}
\usepackage[utf8]{inputenc}
\usepackage{lmodern}
\usepackage{textcomp}
\usepackage[margin=2.5cm]{geometry}
\usepackage{amsmath,amssymb}
\usepackage{xcolor}
\usepackage{tikz}
\usepackage{listings}
\usepackage[colorlinks,allcolors=blue!50!black]{hyperref}
\usepackage{cleveref}
% The layers and the style the TikZ backend of DisCoPy draws with.
\pgfdeclarelayer{edgelayer}
\pgfdeclarelayer{nodelayer}
\pgfsetlayers{edgelayer,nodelayer,main}
\tikzstyle{none}=[inner sep=0mm]
\lstset{
  basicstyle=\ttfamily\footnotesize, columns=fullflexible, keepspaces,
  upquote, showstringspaces=false, backgroundcolor=\color{black!4},
  frame=single, rulecolor=\color{black!4}, literate={→}{{$\to$}}1}
\newenvironment{gap}{\begin{quote}\color{orange!60!black}%
  \textbf{Not written yet.} }{\end{quote}}
